One of the risks in mixed human-AI orchestration is confusing visible activity with durable progress. Many runtime signals are useful for operator awareness: typing indicators, token streams, progress ticks, heartbeats, presence updates, and transient logs can all help users understand that execution is happening. However, those signals are not durable business truth by default. Storing them as if they were equivalent to business events creates noisy history and makes it harder to identify what actually mattered.

In Boreal, this distinction is represented by separating ephemeral transport signals from durable request history. Ephemeral signals may travel through local IPC, browser bridges, streaming transport, or peer channels. Durable history is promoted only when the outcome has business meaning. Examples include a fulfillment state change, a blocker summary, a delivery-ready summary, an artifact reference, or a payment-relevant fact. This keeps the append-only ledger meaningful rather than turning it into a dump of every token and heartbeat.

The distinction is even more important for embodied work. A stream of progress messages does not prove that a site was visited. A verbose model trace does not prove that a human witness observed a handoff. If the system treats runtime chatter as equivalent to completion evidence, it will make generative plan collapse harder to detect. Durable proof has to be attached through the right object family, typically \texttt{Artifact} plus corresponding \texttt{RequestEvent} or \texttt{Fulfillment} transitions.

This separation also improves system design discipline. Interfaces can remain responsive and transparent without inflating the durable history model. Operators can still see transient execution signals, but business state remains tied to explicit promotions. In other words, the model supports visibility without sacrificing auditability.
